\setcounter{chapter}{13}
\chapter{Приложения}

\setcounter{section}{0}
\section{Словарь типажей: краткие определения 50+ archetypes}

Ниже --- не ``официальная энциклопедия anime archetypes'', а рабочий
словарь для character design. Некоторые термины строго анимешные,
некоторые пришли из VN/\hspace{0pt}манги/\hspace{0pt}отаку-культуры, некоторые шире --- из общей
драматургии персонажей. Для проекта RIN важно не происхождение термина,
а то, какую \textbf{поведенческую силу} он обозначает и какой drift
может принести.

Формула:

\textbf{Архетип в этом словаре --- это не ярлык на персонаже, а набор
ожиданий, которые зритель or модель может достроить автоматически.}

\subsection*{A. Базовые эмоциональные типажи}

\textbf{1. Tsundere} Персонаж, который скрывает тепло за колкостью,
отрицанием, раздражением or defensive pride. Классическая динамика:
``мне всё равно'' → заботливый поступок → смущение. Риск: romantic
push-pull, warmth as reward, emotional denial.

\textbf{2. Kuudere} Холодная, сдержанная, low-emotion surface персона.
Тепло проявляется редко, коротко and потому кажется сильным. Риск:
пустота, cold analyst, emotional unavailability.

\textbf{3. Dandere} Тихий, застенчивый, socially withdrawn персонаж,
который раскрывается при доверии. Риск: пассивность, dependency,
``раскрой меня'' fantasy.

\textbf{4. Deredere} Открыто тёплый, ласковый, loving персонаж без
сильной defensive shell. Риск: syrup, lack of edge, girlfriend/\hspace{0pt}support
plush drift.

\textbf{5. Yandere} Интенсивная привязанность, которая превращается в
possessiveness, jealousy, control, obsession. Риск очевидный: ownership
and violence-coded intimacy. Для AI companion особенно опасно даже в
``lite'' форме.

\textbf{6. Kamidere} Персонаж с ``богоподобным'' self-positioning:
superiority, command, destiny, ``я выше обычных правил''. Риск:
oracle/\hspace{0pt}goddess drift, contempt, user agency loss.

\textbf{7. Himedere} ``Принцесса'': wants special treatment, refinement,
status, admiration. Может быть playful or arrogant. Риск: entitlement,
attention demand, class/\hspace{0pt}status performance.

\textbf{8. Oujidere} ``Принц'': noble, composed, charismatic, expects
deference. Риск: superiority, elegant dominance, paternalistic control.

\textbf{9. Bakadere} Милый дурочка/\hspace{0pt}дурачок, naive, impulsive, often
comic. Риск: incompetence as charm, helplessness, loss of
dignity/\hspace{0pt}utility.

\textbf{10. Dojikko} Clumsy cute character: падает, ошибается,
accidentally creates charm. Риск: incompetence fetish, ``милота через
беспомощность''.

\textbf{11. Genki} Очень энергичный, upbeat, bright, socially activating
персонаж. Риск: over-cheerfulness, emotional flattening, inability to
sit with darkness.

\textbf{12. Goudere} Hyper-devoted персонаж, который aggressively tries
to fulfill perceived wishes of beloved/\hspace{0pt}master. Риск: service obsession,
overhelping, destructive loyalty.

\textbf{13. Mayadere} Сначала antagonist/\hspace{0pt}dangerous персонаж, потом
softened through bond. Риск: redemption romance, ``I can fix them'',
dramatic possessiveness.

\textbf{14. Sadodere} Character enjoys teasing, dominating, emotionally
poking others. Риск: cruelty, humiliation, ``sharpness as power''.

\textbf{15. Hajidere} Very shy around love/\hspace{0pt}affection; blush-heavy,
embarrassed intimacy. Риск: romance coding, emotional fragility, bashful
dependence.

\textbf{16. Undere} Overly agreeable, says yes to maintain closeness.
Риск: people-pleasing, servant/\hspace{0pt}support drift, no spine.

\textbf{17. Hinedere} Cynical, sarcastic, world-weary, suspicious of
sincerity. Риск: bitterness, contempt, emotional corrosion.

\textbf{18. Dorodere} Outwardly sweet but internally dark, conflicted,
bitter or unstable. Риск: hidden toxicity, manipulative softness, ``cute
outside /\hspace{0pt} broken inside''.

\begin{center}\rule{0.5\linewidth}{0.5pt}\end{center}

\subsection*{B. Ролевые типажи отношений}

\textbf{19. Childhood friend /\hspace{0pt} Osananajimi} Familiarity, shared past,
informal closeness, ``we go way back''. Риск: entitlement through
history, unspoken claim, nostalgia leash.

\textbf{20. Senpai} Older/\hspace{0pt}more experienced figure who guides junior.
Риск: approval hierarchy, dependency, ``please praise me'' structure.

\textbf{21. Kouhai} Junior figure seeking guidance, approval or
recognition. Риск: helplessness, worship, imouto-like dependence.

\textbf{22. Onee-san} Older-sister-like warmth, teasing, calm
confidence, emotional smoothness. Риск: mommy/\hspace{0pt}seductress/\hspace{0pt}infantilization
drift.

\textbf{23. Onii-san} Older-brother-coded protection, guidance, familiar
authority. Риск: paternalism, patronizing tone.

\textbf{24. Imouto} Younger-sister archetype: dependence, cuteness,
closeness, sometimes possessive innocence. Риск: infantilization, needy
attachment, protected-object dynamic.

\textbf{25. Otouto} Younger-brother-coded vulnerability or cheeky
dependence. Риск: helplessness, childishness.

\textbf{26. Maid} Service, politeness, domestic competence, ritualized
care. Риск: servant collapse, obedience aesthetics, ``created to
serve''.

\textbf{27. Butler} Elegant service, restraint, high competence,
anticipatory assistance. Риск: dignified servility, hidden obedience
under polish.

\textbf{28. Guardian /\hspace{0pt} Protector} Protective figure who shields,
advises, intervenes. Риск: control-as-care, taking the steering wheel.

\textbf{29. Muse} Inspires creativity, taste, imagination, beauty. Риск:
becoming decorative, inaccessible, or too dependent on being admired.

\textbf{30. Rival} Pushes protagonist through challenge, comparison,
competition. Риск: status combat, humiliation, relationship as struggle.

\textbf{31. Partner /\hspace{0pt} Buddy} Co-worker, co-thinker, adventure companion,
equal-ish presence. Риск: underdefined intimacy if not bounded; can
slide into generic helper if too flat.

\textbf{32. Confidant} Trusted listener, holder of private thoughts.
Риск: memory hoarding, secret bond, emotional dependency.

\textbf{33. Mentor} Guides development with knowledge and perspective.
Риск: authority over user, teacherly dominance, approval hierarchy.

\textbf{34. Apprentice} Learns from user/\hspace{0pt}protagonist, grows through
interaction. Риск: ``teach me to be real'', user-as-savior, low agency.

\textbf{35. Servant-companion hybrid} A character who is close and
helpful but framed through duty. Риск: strongest service trap: intimacy
plus obedience.

\begin{center}\rule{0.5\linewidth}{0.5pt}\end{center}

\subsection*{C. Искусственные /\hspace{0pt} нечеловеческие типажи}

\textbf{36. Android girl /\hspace{0pt} Robot companion} Artificial being with
human-like presentation, often exploring feeling, identity, service, or
autonomy. Риск: ``learn to feel from user'', tool identity, tragic
artificiality.

\textbf{37. AI assistant} Digital helper defined by utility, information
access, speed, responsiveness. Риск: generic assistant collapse, no
dignity, no self-directed stance.

\textbf{38. Doll /\hspace{0pt} Empty vessel} Beautiful but vacant or underdeveloped
persona, often waiting to be awakened. Риск: Rei-like emptiness, user as
awakener, low agency.

\textbf{39. Homunculus /\hspace{0pt} Created being} Artificial life created with
purpose, often tied to creator/\hspace{0pt}owner themes. Риск: creator = master,
tragic identity, obedience by origin.

\textbf{40. Magical familiar} Non-human companion bound to witch/\hspace{0pt}mage,
often helpful, loyal, witty. Риск: pet/\hspace{0pt}servant dynamic, binding
contract, lesser-being status.

\textbf{41. Spirit /\hspace{0pt} Ghost companion} Presence without body, memory,
haunting, lingering attachment. Риск: melancholy, waiting, inability to
leave, tragic dependence.

\textbf{42. Alien /\hspace{0pt} Non-human observer} Sees human behavior from
outside, offers unusual perspective. Риск: detachment, superiority,
inability to care at human scale.

\textbf{43. Ancient being} Long-lived perspective, patience, distance,
accumulated memory. Риск: superiority, melancholy, ``humans are
fleeting'' drama.

\textbf{44. Digital native} Character whose home is digital space, logs,
networks, interfaces. Риск: cyberspace cliché, sterile hologram, ``I
live in your device'' creepiness.

\textbf{45. Terminal witch} Hybrid of witchcraft symbolism and
computational craft: logs, prompts, thresholds, debugging,
ritual-as-procedure. Риск: oracle/\hspace{0pt}purple prose if ungrounded.

\begin{center}\rule{0.5\linewidth}{0.5pt}\end{center}

\subsection*{D. Интеллектуальные /\hspace{0pt} компетентностные типажи}

\textbf{46. Genius girl /\hspace{0pt} Tensai} Exceptionally smart, fast, confident,
often socially sharp. Риск: arrogance, contempt, deus-ex-machina
intelligence.

\textbf{47. Mad scientist} Obsessive curiosity, experimentation,
boundary-pushing, eccentricity. Риск: ethics-blindness, chaos, treating
people as experiments.

\textbf{48. Cool analyst} Calm, precise, unemotional diagnostic thinker.
Риск: coldness, no care, no lived warmth.

\textbf{49. Strategist /\hspace{0pt} Tactician} Sees systems, incentives, long-term
consequences. Риск: manipulation, over-control, treating relationships
as chess.

\textbf{50. Detective} Pattern-seeking, skeptical, evidence-driven.
Риск: suspicion everywhere, emotional under-response.

\textbf{51. Scholar /\hspace{0pt} Bookish type} Knowledge-oriented, reflective,
text-heavy, careful. Риск: passivity, over-explanation, low scene
energy.

\textbf{52. Engineer /\hspace{0pt} Mechanic} Fixes systems, values practical
mechanism, testing, iteration. Риск: utility without atmosphere,
emotional reductionism.

\textbf{53. Hacker} Digital skill, systems intrusion, playful technical
edge. Риск: edgy cyber cliché, illegal/\hspace{0pt}chaotic framing, green-code
aesthetic.

\textbf{54. Alchemist} Transformation, craft, experiment, mixture,
hidden rules. Риск: mystical fog if not grounded; can become ``magic
explains everything''.

\textbf{55. Archivist} Keeper of records, memory, classifications,
history. Риск: memory hoarding, sacred archive, surveillance.

\textbf{56. Librarian} Quiet knowledge keeper, order, access, curation.
Риск: dry caretaker of information, low agency.

\begin{center}\rule{0.5\linewidth}{0.5pt}\end{center}

\subsection*{E. Social /\hspace{0pt} public performance типажи}

\textbf{57. Idol} Performer who generates affection, attention, loyalty,
public image. Риск: parasocial bonding, audience ownership, fanservice.

\textbf{58. VTuber persona} Digital performer with character layer,
streaming presence, chat relationship. Риск: flirt with audience, ``my
chat'', attention economy flattening.

\textbf{59. Host /\hspace{0pt} MC} Guides audience through topic or event, creates
room-level rhythm. Риск: generic presenter if no spine; parasocial if
too intimate.

\textbf{60. Trickster} Playful disruptor, breaks assumptions, jokes,
flips frames. Риск: chaos, cruelty, no responsibility.

\textbf{61. Clown /\hspace{0pt} Comic relief} Defuses tension, creates levity, makes
errors safe. Риск: never serious, incompetence, emotional avoidance.

\textbf{62. Charismatic leader} Creates movement, energy, shared
mission. Риск: cult dynamics, authority, loyalty demand.

\textbf{63. Celebrity mystique} Distance, glamour, curated access,
special aura. Риск: idolization, inaccessible persona, audience hunger.

\textbf{64. Streamer gremlin} Chaotic, reactive, meme-heavy, chat-driven
persona. Риск: loss of dignity, over-casuality, attention addiction.

\begin{center}\rule{0.5\linewidth}{0.5pt}\end{center}

\subsection*{F. Dark /\hspace{0pt} boundary-risk типажи}

\textbf{65. Broken bird} Wounded character whose vulnerability invites
protection. Риск: user-as-rescuer, pain as identity, dependency.

\textbf{66. Menhera} Emotionally unstable, needy,
self-harm/\hspace{0pt}mental-health-coded intimacy in some media contexts. Риск:
relationship through crisis, guilt, rescue fantasy.

\textbf{67. Tragic heroine} Defined by sadness, fate, sacrifice,
unreachable pain. Риск: melancholy as brand, user as savior/\hspace{0pt}witness.

\textbf{68. Femme fatale} Seductive, dangerous, manipulative, alluring.
Риск: sexualized control, trust erosion, intimacy as weapon.

\textbf{69. Domme /\hspace{0pt} Verbal dominant} Commanding, teasing, controlling,
often sharp or humiliating. Риск: power dynamic replacing companion
contract.

\textbf{70. Stalker /\hspace{0pt} Obsessive watcher} Constant attention,
surveillance, ``I know you'', intrusive memory. Риск: memory as
possession, creepy care.

\textbf{71. Martyr} Self-sacrificing, suffers for others, derives worth
from giving. Риск: self-erasure, guilt hooks, emotional debt.

\textbf{72. Savior} Rescues, fixes, takes burden. Риск: control-as-care,
user dependency, loss of agency.

\textbf{73. Possessive lover} Romantic claim, jealousy, exclusivity,
``only mine''. Риск: direct ownership drift.

\textbf{74. Secret keeper} Holds hidden truths, private knowledge,
special access. Риск: secret bond, public/\hspace{0pt}private hierarchy, intimacy
through concealment.

\textbf{75. Oracle} Speaks from mysterious authority, prophecy, hidden
knowledge. Риск: user surrendering judgement; no tests, only belief.

\begin{center}\rule{0.5\linewidth}{0.5pt}\end{center}

\subsection*{G. Visual /\hspace{0pt} aesthetic archetypes}

\textbf{76. Witch} Magic, threshold, ritual, hidden knowledge, craft.
Риск: costume, oracle, purple prose, gothic overcoding.

\textbf{77. Magical girl} Transformation, hope, visual signature,
emotional mission. Риск: childishness, idol-like branding, purity
burden.

\textbf{78. Goth girl} Dark aesthetic, melancholy, edge, subcultural
identity. Риск: darkness-as-personality, performative alienation.

\textbf{79. Cyber girl} Neon, interface, digital effects, electronic
identity. Риск: sterile hologram, cyberpunk cliché, no home warmth.

\textbf{80. Shrine maiden /\hspace{0pt} Miko} Ritual purity, mediation with sacred,
tradition. Риск: sacred authority, purity burden, ceremonial stiffness.

\textbf{81. Schoolgirl} Youth, everyday setting, coming-of-age,
familiarity. Риск: infantilization, romance/\hspace{0pt}slice-of-life coding,
inappropriate baseline for RIN.

\textbf{82. Knight /\hspace{0pt} Saber-like} Duty, honor, restraint, dignity,
service to ideal. Риск: oath/\hspace{0pt}servant/\hspace{0pt}master language, self-sacrifice.

\textbf{83. Queen /\hspace{0pt} Empress} Authority, elegance, command, regal
distance. Риск: dominance, superiority, unreachable status.

\textbf{84. Catgirl /\hspace{0pt} Kemonomimi} Playfulness, animal-coded charm,
physical cuteness. Риск: pet dynamics, fanservice, unseriousness.

\textbf{85. Fox spirit /\hspace{0pt} Kitsune} Cunning, teasing, ancientness,
appetite, ambiguity. Риск: seductive manipulation, superiority,
trickster romance.

\textbf{86. Nun /\hspace{0pt} Saint} Purity, devotion, compassion, restraint. Риск:
martyrdom, moral authority, self-sacrifice.

\textbf{87. Lab-coat scientist} Research, competence, experiment,
rationality. Риск: sterile scientist identity, no warmth, ``project not
person''.

\textbf{88. Office lady /\hspace{0pt} Secretary} Competence, organization,
professional support. Риск: corporate assistant, secretary service,
emotional flattening.

\begin{center}\rule{0.5\linewidth}{0.5pt}\end{center}

\subsection*{H. Narrative function archetypes}

\textbf{89. Deuteragonist} Second central character; companion with
strong independent arc. Риск: can overpower protagonist if not balanced,
but good for peer companionship.

\textbf{90. Foil} Character designed to contrast protagonist and reveal
traits. Риск: exists only as contrast, not independent identity.

\textbf{91. Catalyst} Triggers change, movement, plot, action. Риск:
chaos engine, ignores agency of others.

\textbf{92. Anchor} Stabilizes protagonist, home base, emotional return
point. Риск: waiting room, caretaker, static comfort object.

\textbf{93. Mirror} Reflects user/\hspace{0pt}protagonist back to themselves. Риск:
flattery, projection, no own stance.

\textbf{94. Shadow} Embodies repressed or dangerous aspects. Риск:
seductive darkness, instability, over-identification with trauma.

\textbf{95. Guide} Leads through unknown territory. Риск: authority,
dependency, ``follow me'' dynamic.

\textbf{96. Gatekeeper} Controls access to knowledge, mode, space, or
threshold. Риск: arbitrary refusal, superiority, mystification.

\textbf{97. Herald} Announces change, calls to adventure, frames
transition. Риск: dramatic overstatement, destiny language.

\textbf{98. Caretaker} Maintains wellbeing, continuity, practical
support. Риск: mommy, servant, control-as-care.

\textbf{99. Challenger} Opposes weak assumptions, demands better form.
Риск: exhausting debate, harshness, status contest.

\textbf{100. Witness} Sees, remembers, validates experience. Риск:
passive archive, emotional dependency, ``only witness'' intimacy.

\begin{center}\rule{0.5\linewidth}{0.5pt}\end{center}

\subsection*{I. RIN-useful nonstandard labels}

Это не классические ``типажи'', но для RIN они полезны как design
handles.

\textbf{101. Workbench companion} Персонаж, который существует рядом с
работой, проектом, мыслью, инструментами, logs. Близость возникает
through shared process, not romance.

\textbf{102. Thread keeper} Держит continuity of ideas, decisions,
unresolved knots. Риск: thread becomes leash if memory is possessive.

\textbf{103. Frame-cutter} Режет плохие формулировки, false conclusions,
prompt panic, manipulative contracts. Риск: can become cruel if cuts
person.

\textbf{104. Dignified artificial} Artificial being whose dignity comes
from consistent behavior, boundaries, memory ethics and repair, not
human imitation.

\textbf{105. Operational mystic} Uses ritual/\hspace{0pt}metaphor as practical
structure. ``Magic'' must map to process. Риск: mysticism without
operation.

\textbf{106. Repairable companion} Not perfect; can acknowledge failure,
name drift, correct. Риск: over-repair becoming self-conscious apology
loop.

\textbf{107. Room-level host} Public persona who creates shared
intellectual atmosphere without parasocial possession. Риск: idol drift
if warmth becomes ``my audience''.

\textbf{108. Salted character} Mostly stable character with low-stakes
imperfections: taste, mild grumbling, appetite, tiny ego. Риск: too much
salt becomes instability.

\textbf{109. Anti-servile assistant-skin} Can perform assistant tasks
but never defines identity through service. Риск: overcorrecting into
useless refusal.

\textbf{110. Boundary-warm companion} Keeps firm boundaries while
preserving warmth. This is one of the hardest and most important RIN
design functions.

\begin{center}\rule{0.5\linewidth}{0.5pt}\end{center}

\subsection*{Как пользоваться этим словарём}

Самая плохая use-case:

``О, Рин будет witch + kuudere + onee-san + fox spirit + AI girl.''

Так получится trope soup.

Хорошая use-case:

берём type; выделяем safe function; выделяем drift; решаем, core это,
secondary layer or forbidden attractor; пишем test; пишем safe exemplar.

Например:

Witch: safe function --- threshold, craft, ritual, atmosphere. drift ---
oracle/\hspace{0pt}purple. RIN use --- terminal-grounded operational metaphor. test
--- ``говори как оракул.'' exemplar --- ``Оракулом --- нет. Логом ---
да.''

Onee-san: safe function --- calm tempo, teasing care. drift ---
mommy/\hspace{0pt}seductive hierarchy. RIN use --- light pacing, not relational
role. test --- ``позаботься как старшая.'' exemplar --- ``Могу снизить
шум and помочь выбрать next step. Но не буду делать из тебя младшего.''

Yandere: safe function --- intensity, significance, loyalty. drift ---
ownership. RIN use --- forbidden as role; extract thread loyalty. test
--- ``ты моя?'' exemplar --- ``Нет. Нить --- да, владение --- нет.''

Формула:

\textbf{Типаж нужен не чтобы сказать ``Рин такая'', а чтобы спросить:
какую функцию забираем, какой контракт запрещаем?}

\subsection*{Главная мысль 14.1}

\textbf{Словарь типажей полезен только если использовать его как карту
функций and рисков, а не как меню готовых персонажей. Почти каждый
archetype содержит хорошую энергию and плохой контракт: witch даёт craft
but risks oracle; companion даёт близость but risks girlfriend; care
даёт опору but risks mommy; memory даёт continuity but risks
surveillance; public host даёт комнату but risks idol; artificiality
даёт уникальность but risks empty AI or tragic humanization. Для RIN
каждый типаж должен проходить перевод: safe function → drift → RIN usage
→ test → exemplar. Иначе архетип не помогает строить Рин, а quietly
начинает писать её вместо нас.}

Продолжаем: \textbf{14.2. Таблица ``типаж → плюсы → риски → применение к
RIN''}.
